\section{Results}
\label{sec:results}

The results address the three questions posed in the Introduction. The controlled mapping and the efficiency--power analyses show how a higher-variance relative feature can still inform. The sports landscape and the quadrant exemplars show why relativisation helps for some KPIs and harms others. The supporting domains show whether the same $(\kappa,\rho)$ geometry appears outside sport.

\subsection{Theoretical Foundation: PEF--Information Content}

The theoretical relationship between $\eta$ and information content (\cref{eq:mi_closed}) was examined in four named scenarios from the theory-aligned idealised probit simulation (\cref{fig:si_idealised_stratified}; Supplementary~\cref{sec:si_probit}) under assumptions (A1)--(A2), with $n=5{,}000$ per cell, 50 Monte Carlo trials, and five-fold logistic cross-validation matching the empirical protocol. The simulation's data-generating process, estimands, and parameter-recovery checks are specified in full in Supplementary~\cref{sec:si_probit}.

\begin{table}[t]
  \centering
  \caption{Scenario-level confirmation from the idealised probit simulation (A1)--(A2).}
  \label{tab:scenarios}
  \begin{tabular}{@{}lccccc@{}}
    \toprule
    \textbf{Scenario} & $\kappa$ & $\rho$ & $\eta$ & $I(X;Y)$ (bits) & ML impr.$^{\dagger}$ \\
    \midrule
    High competitive dynamics      & 2.0  & $-0.3$  & 0.780 & 0.029 & 9.4\% \\
    Moderate competitive dynamics  & 1.2  & $-0.15$ & 0.870 & 0.044 & 8.2\% \\
    Environmental dynamics         & 1.1  & 0.4     & 1.665 & 0.087 & 8.6\% \\
    Balanced (Fisher baseline)     & 1.0  & 0.0     & 1.000 & 0.056 & 7.9\% \\
    \bottomrule
  \end{tabular}

  \medskip
  \begin{minipage}{0.95\linewidth}\small
    $^{\dagger}$Mean five-fold cross-validated accuracy improvement (relative minus absolute features) across 50 Monte Carlo trials per scenario cell ($\delta/\sigma_{\mathrm{A}}=1$). Standard errors are $\approx$0.8--1.0 percentage points (Supplementary~\cref{sec:si_probit}).
  \end{minipage}
\end{table}

On the admissible $(\kappa,\rho)$ grid at fixed $\delta/\sigma_{\mathrm{A}}=1$, analytic $\eta$ and $I(X;Y)$ from \cref{eq:mi_closed} are strongly correlated ($r\approx 0.87$ across 16 grid cells in the idealised simulation). The four scenarios in \cref{tab:scenarios} span the quadrant space under controlled (A1)--(A2) conditions; the first two illustrate the efficiency--power tension ($\eta<1$ with positive mean ML improvement), consistent with \cref{sec:eff_power_theory}.

\subsection{Sports KPI Landscape}

This inventory is the empirical setting of the help-versus-harm question. The full sports inventory spans all four quadrants (Supplementary~\cref{fig:si_kpi_labelled}; $\PEFtotalStudies{}$ KPI studies, \cref{sec:outcome_defs}). Rugby union ({URC}) contributes $\PEFrugbyNkpis{}$ KPIs and football (English Championship) $\PEFfootNkpis{}$; per-KPI $(\hat\kappa,\hat\rho,\hat\eta)$ estimates and season-to-season drift appear in Supplementary~\cref{fig:si_kpi_labelled,fig:si_ipred_vs_dml}. Domain-level landscape summaries are collected in \cref{tab:validation}: rugby mean $\hat\eta=\PEFrugbyMeanEta$ (95\% CI: $[\PEFrugbyCIlo,\PEFrugbyCIhi]$; $\PEFrugbySuccess\%$ of KPIs with $\hat\eta>1$), football mean $\hat\eta=\PEFfootMeanEta$ (95\% CI: $[\PEFfootCIlo,\PEFfootCIhi]$; $\PEFfootSuccess\%$ with $\hat\eta>1$). These figures characterise the empirical $(\kappa,\rho)$ distribution, notably the Q3-heavy structure of invasion-game KPIs ($\PEFlandscapeQthreeN{}$ of $\PEFtotalStudies{}$ studies; $\PEFlandscapeQthreePct\%$). They do not serve as the primary confirmatory test.

KPIs with positive correlation (e.g., kicks from hand, carries) concentrate in Quadrants~1--2 and typically yield $\hat\eta>1$; KPIs with negative or near-zero correlation (e.g., turnovers won, penalties conceded) fall predominantly in Quadrants~3--4 with $\hat\eta\le 1$, yet several still show positive ML improvement in the full inventory (Supplementary~\cref{fig:si_ipred_vs_dml}), consistent with the efficiency--power tension (\cref{sec:eff_power_theory}). Mechanistic confirmation of the taxonomy's directional predictions is assessed through the four quadrant exemplars below (\cref{fig:pef_ml,sec:exemplars}), not through pooled success rates over the Q3-dominated inventory.

\subsection{Efficiency--Power Alignment Across Quadrants}
\label{sec:eff_power}

This section addresses how a less efficient relative feature can still inform. The four quadrants differ not only in mean $\eta$ but in the relationship between efficiency and predictive gain. Three supplementary analyses characterise this structure.

\textbf{Idealised simulation under (A1)--(A2).} On the theory-aligned probit grid (Supplementary~\cref{fig:si_idealised_stratified}; $n=5{,}000$, 50 trials, five-fold CV), all 16 out of 16 Quadrant~4 cells with $\eta<1$ produced positive mean ML improvement. At fixed $\delta/\sigma_{\mathrm{A}}$, $\eta$ and $I(X;Y)$ are strongly correlated (within-slice $r\approx0.87$ at the empirical median $\delta/\sigma_{\mathrm{A}}=\PEFmedianDeltaRatio$). However, the pooled-across-$\delta/\sigma_{\mathrm{A}}$ correlation is negative ($r=-0.70$), because signal strength shifts the $I(X;Y)$ level independently of the $(\kappa,\rho)$ position. The practical consequence is that iso-$\eta$ and iso-$I(X;Y)$ contours diverge once $\delta/\sigma_{\mathrm{A}}$ varies across KPIs (Supplementary~\cref{fig:si_iso_eta_I}).

\textbf{Quadrant distribution.} Of the \PEFtotalStudies{} sports KPI studies, $\PEFlandscapeQthreeN{}$ ($\PEFlandscapeQthreePct\%$) lie in Q3: the negative-correlation regime where competitive dynamics dominate and absolute features are typically preferred. This reflects a structural feature of professional invasion-game sport: high-volume counts (passes, pressures, duels) are anti-correlated across opponents. Because the dataset is Q3-heavy, any aggregate statistic pooled across all KPIs primarily characterises this dominant regime rather than the mechanism across the full quadrant space (Supplementary~\cref{tab:si_quad_landscape}). The exemplar approach (\cref{tab:exemplars}) directly addresses this by selecting one KPI per quadrant, allowing the taxonomy's directional predictions to be assessed without the distributional confound.

\textbf{Quadrant~4 and the Bayes bound.} Team-stratified bootstrap intervals for tackles (rugby, Q4) lie entirely below $\eta=1$ (Supplementary~\cref{sec:si_qc}). Under (A1)--(A2) with equal priors, the Gaussian Bayes accuracy of the absolute feature is the theoretical comparator. For the Q4 confirmatory KPI, team-blocked CV accuracies for absolute and relative features are close to each other and sit below that bound. The gap is expected: real match outcomes are not generated from a single KPI by the probit link~(A2). Relativisation does not lift accuracy here, matching the Q4 row of \cref{tab:exemplars}. This is the limited-signal Quadrant~4 case, not the efficiency--power tension. Relative gain despite $\eta<1$ appears in the simulation grid and in parts of the landscape (\cref{fig:si_idealised_stratified,fig:si_ipred_vs_dml}). Statistical efficiency and information content remain related but distinct: the value of relativisation cannot be read from $\eta$ alone (\cref{sec:eff_power_theory}).

\subsection{Quadrant Exemplars and Signal-Strength Context}
\label{sec:exemplars}

The confirmatory test of help versus harm is the sign of $\Delta\mathrm{ML}$ at matched signal. \Cref{fig:pef_ml} shows the four confirmatory exemplars on the $\hat\eta$--$\Delta\mathrm{ML}$ plane, with rucks won as a high-$\eta$ foil at negligible $\delta/\sigma_{\mathrm{A}}$ (\cref{tab:exemplars}). The full inventory remains a landscape characterisation (Supplementary~\cref{fig:si_kpi_labelled,fig:si_ipred_vs_dml}).

Each $\Delta\mathrm{ML}$ in \cref{tab:exemplars} comes from \emph{univariate} logistic regression (one KPI feature per model: $X_{\mathrm{A}}$ or $X_{\mathrm{A}}-X_{\mathrm{B}}$). Percentage changes are modest in absolute terms yet substantively informative. The confirmatory question is the \emph{sign}. $\Delta\mathrm{ML}>0$ means the relative feature classifies match outcome more accurately than $X_{\mathrm{A}}$ alone. $\Delta\mathrm{ML}<0$ means relativisation is predictively harmful: the difference is a worse univariate classifier than the absolute KPI. That harm is the empirical content of the Quadrant~3 recommendation, not a failed gain. Statistical harm ($\eta<1$) and predictive harm ($\Delta\mathrm{ML}<0$) can disagree. Q3 combines both. Rucks won has $\eta>1$ with $\Delta\mathrm{ML}<0$. Q4 has $\eta<1$ with near-zero $\Delta\mathrm{ML}$ (\cref{fig:pef_ml,tab:exemplars}). The present design does not construct an optimal multi-feature match predictor. Prior rugby and football work that motivates this study combined many KPIs in a single classifier, where per-feature changes can compound, cancel, or interact (\cref{sec:discussion}). The single-KPI protocol is conservative by construction. Team-blocked cross-validation (\cref{sec:ml_cv}) avoids the optimistic bias of random match folds when the same side recurs across fixtures. It also reduces leakage of a team's early and late home matches across train and test partitions; season-blocked validation is noted as a further tightening step in \cref{sec:discussion}.

The four confirmatory KPIs occupy distinct regions of the $(\kappa,\rho)$ plane (\cref{fig:pef_landscape,tab:exemplars}) at comparable signal strength; they are not interchangeable summaries of their quadrants. Because $\eta$ and $I(X;Y)$ respond nonlinearly to $(\kappa,\rho)$ and to $\delta/\sigma_{\mathrm{A}}$ (\cref{eq:mi_closed}), each exemplar illustrates \emph{local} regime behaviour. Interpreting a quadrant label as a uniform prescription would therefore overstate what the taxonomy claims.

\begin{table}[t]
  \centering
  \caption{Quadrant exemplars: four KPIs spanning the $(\kappa,\rho)$ parameter space.}
  \label{tab:exemplars}
  \begin{tabular}{@{}llccccc@{}}
    \toprule
    \textbf{Q} & \textbf{KPI (sport)} & $\kappa$ & $\rho$ & $\eta$ & $\delta/\sigma_{\mathrm{A}}$ & $\Delta\mathrm{ML}$ (\%) \\
    \midrule
    Q1 & Kick metres (rugby)              & $\PEFexQoneKappa$   & $\PEFexQoneRho$   & $\PEFexQoneEta$   & $\PEFexQoneDeltaRatio$   & $\PEFexQoneDml$ \\
    Q2 & Long balls (football)            & $\PEFexQtwoKappa$   & $\PEFexQtwoRho$   & $\PEFexQtwoEta$   & $\PEFexQtwoDeltaRatio$   & $\PEFexQtwoDml$ \\
    Q3 & Passes (football)                & $\PEFexQthreeKappa$ & $\PEFexQthreeRho$ & $\PEFexQthreeEta$ & $\PEFexQthreeDeltaRatio$ & $\PEFexQthreeDml$ \\
    Q4 & Goalkeeper long balls (football) & $\PEFexQfourKappa$  & $\PEFexQfourRho$  & $\PEFexQfourEta$  & $\PEFexQfourDeltaRatio$  & $\PEFexQfourDml$ \\
    \bottomrule
  \end{tabular}

  \medskip
  \begin{minipage}{0.95\linewidth}\small
    Two-season pooled estimates ($\kappa$, $\rho$, $\eta$, $\delta/\sigma_{\mathrm{A}}$) and team-blocked five-fold cross-validated logistic change ($\Delta\mathrm{ML}$; \cref{sec:ml_cv}). Signal strengths ($\delta/\sigma_{\mathrm{A}}\approx\PEFexDeltaRatioLo$--$\PEFexDeltaRatioHi$) are broadly comparable across exemplars. The table tests sign at matched signal: Q1 and Q2, $\Delta\mathrm{ML}>0$ with $\eta>1$; Q3, predictive harm ($\Delta\mathrm{ML}<0$) with $\eta<1$; Q4, $\eta<1$ with near-zero $\Delta\mathrm{ML}$. Relative gain despite $\eta<1$ is shown in the idealised grid and the full inventory (\cref{sec:eff_power,fig:si_ipred_vs_dml}), not by a large Q4 exemplar gain.
  \end{minipage}
\end{table}

The Q4 exemplar (goalkeeper long balls, $\eta=\PEFexQfourEta$) has $\Delta\mathrm{ML}=\PEFexQfourDml\%$ at team-blocked cross-validation. The change is near zero at this $\delta/\sigma_{\mathrm{A}}$, as expected when pairing inflates variance and signal is limited. This row is not a demonstration of relative benefit despite $\eta<1$. Cases with $\eta<1$ and $\Delta\mathrm{ML}>0$ appear in the full inventory (\cref{fig:si_ipred_vs_dml,sec:eff_power}). Here $\kappa\approx\PEFexQfourKappa$, so $\eta<1$ is driven by negative correlation ($\rho=\PEFexQfourRho$) rather than variance asymmetry. The KPI therefore sits on the boundary between Quadrants~3 and~4.

The Q3 exemplar (passes, football; $\rho=\PEFexQthreeRho$, $\eta=\PEFexQthreeEta$) is the complementary result. Strong negative correlation places this KPI deep in the low-efficiency regime, and $\Delta\mathrm{ML}=\PEFexQthreeDml\%$ shows that relativisation is predictively harmful. The absolute feature is the better univariate classifier, which is the Quadrant~3 recommendation in \cref{tab:quadrants}. Under (A1)--(A2) the idealised probit grid cannot produce $\Delta\mathrm{ML}<0$ (Supplementary~\cref{sec:si_probit}). With a Bayes-sufficient relative feature, pairing does not reduce accuracy. Empirical harm is therefore a property of real match outcomes: the difference can discard useful absolute-level information, or it can feed the classifier a noisier feature. For high-volume counts such as passes, the difference can discard game volume or tempo.

Rucks won (rugby; Q2: $\kappa=\PEFexRucksKappa$, $\rho=\PEFexRucksRho$, $\eta=\PEFexRucksEta$, $\delta/\sigma_{\mathrm{A}}=\PEFexRucksDeltaRatio$) shows that harm is not confined to $\eta<1$. Despite the highest $\eta$ in the dataset, $\Delta\mathrm{ML}=\PEFexRucksDml\%$ because the two teams are almost equally matched on average. Pairing efficiency without signal still yields a worse predictor than the absolute KPI.

\subsection{Supporting Validation: Cross-Domain Results}

The third question is whether similar $(\kappa,\rho)$ geometry appears beyond sport. The supporting domains are second-moment checks, not match-outcome replications. \Cref{tab:validation} summarises results across all six domains. They show how empirical $(\kappa,\rho)$ structures propagate into $\hat\eta$ and landscape efficiency rates (proportion with $\hat\eta>1$; \cref{sec:outcome_defs}):
\begin{itemize}
  \item \textbf{Healthcare} \parencite{nhanes2017pbxo}: pooled $\eta=\PEFhealthMeanEta$ (bootstrap 95\% CI: $[2.448,2.622]$) across $N=10{,}352$ {NHANES} participants pairing systolic and diastolic blood pressure from the same oscillometric measurement (Quadrant~2). Pulse pressure ($\mathrm{SBP}-\mathrm{DBP}$) is the clinically established relative feature. The one healthcare row in \cref{tab:validation} is that single paired study, not a participant count.
  \item \textbf{Clinical genomics} \parencite{brunner2014earlybreast,geoGSE47462}: mean $\eta=\PEFgenomicsMeanEta$ ({SD} $=\PEFgenomicsSdEta$) across $\PEFgenomicsN$ genes ranked by paired-differential signal in {GSE47462}; $\PEFgenomicsSuccess\%$ of genes yielded $\eta>1$.
  \item \textbf{Finance}: mean $\eta=\PEFfinanceMeanEta$ ({SD} $=\PEFfinanceSdEta$) across $\PEFfinanceN$ {S\&P}~100 stocks relativised to daily {S\&P}~500 returns (2020--2023 pinned snapshot); $\PEFfinanceSuccess\%$ with $\hat\eta>1$.
  \item \textbf{Manufacturing} \parencite{openml752bosch}: mean $\eta=\PEFmfgMeanEta$ ({SD} $=\PEFmfgSdEta$) on {Bosch} {OpenML} 752 {CNC} cycles ($N=\PEFmfgN$ axis-wise pairings). Approximately $\PEFmfgSuccess\%$ of pairings yielded $\eta>1$.
\end{itemize}
Summary statistics for all six domains are collected in \cref{tab:validation}; the sports landscape with supporting-domain overlay markers appears in \cref{fig:pef_landscape}. Detailed per-study breakdowns for the non-sports domains are not separately visualised given the cross-sectional nature of those datasets.
